\documentclass[conference]{IEEEtran}
\IEEEoverridecommandlockouts
% The preceding line is only needed to identify funding in the first footnote.
% If that is unneeded, please comment it out.
\usepackage{cite}
\usepackage{amsmath,amssymb,amsfonts}
\usepackage{algorithmic}
\usepackage{graphicx}
\usepackage{textcomp}
\usepackage{xcolor}
\usepackage{float}
\usepackage{subcaption}
\usepackage{amsmath}
\usepackage{algpseudocode}
\usepackage{algorithm}

% ====== 新增以下這段來支援中文 ======
\usepackage[AutoFakeBold=true, AutoFakeSlant=true]{xeCJK}
% 設定主要的繁體中文字體，您可以根據您的作業系統或編譯環境更改
\setCJKmainfont{標楷體}   % 如果想要比較正式的襯線字體

\def\BibTeX{{\rm B\kern-.05em{\sc i\kern-.025em b}\kern-.08em
    T\kern-.1667em\lower.7ex\hbox{E}\kern-.125emX}}
\begin{document}

\title{CUDA-Accelerated Implementation of Genetic Algorithm-Optimized Ant Colony Optimization for Pipe Routing \\
{\footnotesize 基於基因演算法最佳化之蟻群演算法於管線佈局問題的 CUDA 加速實現}
}

\author{\IEEEauthorblockN{1\textsuperscript{st} Hung Ming-Chun}
\IEEEauthorblockA{\textit{Department of Computer Science and} \\
\textit{Information Engineering}\\
\textit{National Taiwan Normal University}\\
Taipei, Taiwan(R.O.C) \\
61447026s@ntnu.edu.tw}
\and
\IEEEauthorblockN{2\textsuperscript{nd} Hsieh Yao-yu}
\IEEEauthorblockA{\textit{Graduate Institute of} \\
\textit{Science Education}\\
\textit{National Taiwan Normal University}\\
Taipei, Taiwan(R.O.C) \\
61345001s@ntnu.edu.tw}
}

\maketitle

\begin{abstract}
% 請在此處填寫摘要
\end{abstract}

\section{Introduction}
隨著現代工程與建築設計的複雜度日益提升，管線佈局（Pipe Routing）
問題已成為設施配置中不可或缺且極具挑戰性的一環。在有限的空間內，
如何有效且安全地佈置多種不同屬性的管線（如瓦斯管、電管、水管等），
同時滿足物理限制與工程規範，是一個典型的 NP-hard 最佳化問題 \cite{b1}。
傳統的佈局方法多依賴工程師的經驗，不僅耗時費力，且難以在複雜環境中保證解的最佳性。\\

為了解決此問題，啟發式演算法（Metaheuristics）逐漸被引入管線佈局的領域。
其中，蟻群演算法（Ant Colony Optimization, ACO）因其優異的正向回饋機制與探索路徑能力，
在陳俊仁與翁世光教授的研究中可以看到蟻群演算法的潛力\cite{b2}。
然而，ACO 演算法的效能高度依賴於超參數
（如費洛蒙權重 $\alpha$、啟發式權重 $\beta$ 等）的設定，
參數選擇不當極易導致演算法陷入局部最佳解或收斂過於緩慢。\\

有鑑於此，本研究提出一種混合式的最佳化架構（GA-ACO），
引入基因演算法（Genetic Algorithm, GA）作為高階參數調校機制，
動態且自動地搜尋 ACO 的最佳參數組合。此外，針對多條管線佈局與 GA 龐大評估帶來的龐大計算負荷，
本研究利用 CUDA 架構對 ACO 的路徑探索過程進行大規模平行化加速（Parallelization），
並針對管線間的安全距離（Clearance）設計了軟性懲罰機制（Soft Penalty），
以確保管線在受限空間中的合法性與合理性。最後，
本研究開發了一套基於 OpenCV 與 Gnuplot 的即時動態視覺化系統，
不僅能直觀地展示管線收斂結果，亦能即時追蹤 GA 的演化過程，為管線佈局問題提供了一個高效、
穩定且具備高度擴展性的解決方案。

\section{Related Work}

\subsection{Pipe Routing Problem}
管線佈局問題（Pipe Routing Problem）廣泛存在於船舶設計、化工廠、建築機電（MEP）與積體電路設計等領域 \cite{b1}。
其核心目標是在給定的三維或二維空間中，為多個起點與終點尋找無碰撞的連接路徑，
同時將總長度、彎折次數及與其他管線的安全距離最小化。早期的研究多採用確定性演算法，
如 Dijkstra 演算法 \cite{b3} 或 A* 演算法 \cite{b4}，雖然能確保找到單一管線的最短路徑，
但在處理多管線且需考慮複雜避讓規則時，其運算時間會隨問題規模呈指數增長。
近年來，研究焦點逐漸轉向啟發式演算法，以期在合理的時間內求得近似最佳解。

\subsection{Genetic Algorithm}
基因演算法（Genetic Algorithm, GA）是一種模擬自然界物種演化過程的群體演算法 \cite{b5}。
透過選擇（Selection）、交配（Crossover）與突變（Mutation）等操作，
GA 能夠在廣大的搜尋空間中進行全域搜索。在工程最佳化問題中，
GA 常被用來解決離散或連續變數的組合問題。在本研究中，GA 並非直接用於規劃管線路徑，
而是被應用於超參數最佳化（Hyperparameter Tuning），
藉由評估不同參數組合（如 $\alpha$, $\beta$, $\rho$, $Q$ 等）所產生的路徑品質，
引導整體架構朝向更穩健的設定演化，克服人工調參的盲目性與耗時性。

\subsection{Ant Colony Optimization}
蟻群演算法（Ant Colony Optimization, ACO）是由 Marco Dorigo 等人提出 \cite{b6}，
模擬自然界螞蟻覓食行為的演算法。螞蟻在行走路徑上會留下費洛蒙（Pheromone），
後續的螞蟻會根據費洛蒙濃度與啟發式資訊選擇路徑，進而促使整個蟻群逐漸收斂至最短路徑。
在管線佈局應用中，ACO 能夠輕易整合各種空間限制與工程懲罰（如轉彎成本），
具備良好的空間適應能力。然而，傳統 ACO 面臨著容易早熟收斂（Premature Convergence）
與參數敏感度高的缺陷，這也是本研究結合 GA 進行優化的主要動機。

\subsection{CUDA Acceleration}
統一計算架構（Compute Unified Device Architecture, CUDA）是 NVIDIA 推出的平行運算平台與編程模型，
允許開發者利用圖形處理器（GPU）強大的多核心運算能力來處理通用計算任務。在群體智慧演算法中，個體
（如螞蟻或染色體）之間的行為多具備高度的獨立性，非常適合映射至 GPU 的平行執行緒（Threads）上。
近年相關文獻指出，將 ACO 移植至 CUDA 架構可大幅縮短演算法的執行時間 \cite{b7}。
本研究進一步利用 CUDA Streams 技術，實現了「基因層級」與「螞蟻層級」的雙重平行化，
極大地提升了 GA-ACO 混合演算法在複雜環境下的執行效率。

%%%%%%%%%%%%%%%%%%%%%%%%%%%%%%%%%%%%%%%%%%%%%%%%%%
\section{Implement Details}

\subsection{Genetic Algorithm}
本研究的 GA 設計旨在為 ACO 尋找最佳的超參數組合。
我們將 ACO 的關鍵參數定義為染色體（Chromosome），
包含費洛蒙權重 $\alpha$、啟發式權重 $\beta$、費洛蒙揮發率 $\rho$、釋放常數 $Q$、
轉彎懲罰權重 $turn\_w$、安全距離懲罰權重 $clear\_w$ 以及擴散半徑 $diffusion\_rad$。
演化過程如下：首先隨機初始化一個大小為 $N$ 的族群，並為了避免「冷啟動死亡（Cold Start Failure）」，
在初始族群中植入一組已知的保底參數。每一代的評估階段，GA 會呼叫 ACO 
進行模擬並計算適應度（Fitness），適應度定義為該組參數所產生之最佳管線組合的總成本
（包含路徑長度、轉彎懲罰與安全距離軟性懲罰）。接著，利用排序法保留前 50\% 的菁英個體，
並對剩餘個體進行算術交配與 20\% 機率的隨機突變，以產生下一代子代。

\subsubsection{Crossover Function and Mutation}
在子代的生成過程中，本研究採用算術交配（Arithmetic Crossover）與均勻突變（Uniform Mutation）。
假設從菁英群體中透過輪盤或隨機挑選出兩組父代個體 $P_1$ 與 $P_2$，其產生的子代 $C$ 基因計算方式為兩者之算術平均：
\begin{equation}
C_{gene} = \frac{P_{1,gene} + P_{2,gene}}{2}
\end{equation}
為維持族群的多樣性並避免陷入局部最佳解，每個生成的子代基因皆有 $p_{m} = 0.2$ 的機率發生突變。突變操作會無視父代特徵，直接於該參數的合理物理邊界 $[L_{bound}, U_{bound}]$ 內，重新生成一均勻分布之隨機數。

\begin{algorithm}[H]
\caption{Genetic Algorithm for ACO Hyperparameter Tuning}
\begin{algorithmic}[1]
\STATE Initialize Population $P$ with size $N$
\STATE Insert known robust parameters into $P[0]$ (Seeding)
\FOR{$gen = 1$ to $GA\_GENERATIONS$}
    \STATE \textbf{Parallel Evaluate:} $Fitness(P_{i}) \leftarrow \text{run\_aco}(P_{i}) \quad \forall i \in \{1..N\}$
    \STATE Sort $P$ in ascending order based on $Fitness$
    \STATE Keep top 50\% of $P$ as Elites
    \FOR{$i = N/2$ to $N-1$}
        \STATE Select Parents $P_a, P_b$ from Elites randomly
        \STATE $P_{i} \leftarrow \text{Crossover}(P_a, P_b)$
        \IF{$rand() < 0.2$}
            \STATE $P_{i} \leftarrow \text{Mutate}(P_{i})$
            \STATE $P_{i}.Fitness \leftarrow -1$ (Require re-evaluation)
        \ENDIF
    \ENDFOR
\ENDFOR
\STATE \textbf{Return} Best Chromosome $P[0]$
\end{algorithmic}
\end{algorithm}


\subsection{Ant Colony Optimization}
在 ACO 的實作中，我們將每種屬性的管線（如瓦斯管與電管）指派給特定類型的螞蟻。
為了支援多起點對多終點的佈局，螞蟻的生成採用餘數與商數的交錯分配法，確保所有起迄組合均被均勻探索。
螞蟻的移動決策基於輪盤賭選擇法（Roulette Wheel Selection），
並引入了 $\epsilon$-Greedy 策略，賦予螞蟻 3\% 的機率進行盲目隨機探索，
以利跳脫局部最佳解。在費洛蒙更新機制上，我們採用了混合釋放策略：
不僅該回合的區域菁英（Iteration Best）會釋放費洛蒙以維持探索動能，
全域歷史最佳解（Global Best）亦會釋放大量費洛蒙以強化主幹道的收斂。
此外，當不同管線發生重疊或違反安全距離（Safe Distance）時，
系統會計算重疊次數並給予極重的軟性懲罰，引導演算法尋找無碰撞的完美配置。

\subsubsection{Cost Function}
為了全面評估管線配置的優劣，本研究設計之成本函數（Cost Function）結合了路徑長度、幾何彎折度以及管線間的實體干涉懲罰。單一螞蟻的路徑成本 $S_{path}$ 定義如下：
\begin{equation}
S_{path} = L + (W_{turn} \times C_{turn}) + (W_{clear} \times P_{clear})
\end{equation}
其中，$L$ 為路徑網格總長度，$C_{turn}$ 為路徑彎折次數，$P_{clear}$ 為進入其他管線安全半徑內的嚴重程度累加值；$W_{turn}$ 與 $W_{clear}$ 則為受 GA 調控的基因權重。
而在全域聯合評估時，若瓦斯管與電管發生實體重疊，我們不採用硬性淘汰（Hard Constraint），而是給予軟性梯度懲罰（Soft Penalty）：
\begin{equation}
S_{global} = \sum S_{path} + (1000 \times C_{overlap})
\end{equation}
此機制允許 GA 在初期接受包含輕微重疊的瑕疵解，並透過梯度引導族群逐步收斂至重疊數 $C_{overlap} = 0$ 的合法佈局。

\begin{algorithm}[H]
\caption{Parallel ACO Routing with Soft Penalty}
\begin{algorithmic}[1]
\FOR{$iter = 1$ to $MAX\_ITER$}
    \STATE Move all ants concurrently in GPU Kernel
    \STATE Evaporate Pheromones: $\tau \leftarrow \tau \times (1 - \rho)$
    \FOR{each pipe type and endpoint}
        \STATE Evaluate $S_{path}$ using Eq (2)
        \STATE Update Iteration Best ant
    \ENDFOR
    \IF{all endpoints reached}
        \STATE Calculate $C_{overlap}$ between different pipe types
        \STATE Evaluate combined score $S_{global}$ using Eq (3)
        \IF{$S_{global} < Best\_S_{global}$}
            \STATE Update Global Best Paths
        \ENDIF
    \ENDIF
    \STATE Deposit Pheromones (Iteration Best $\times 30\%$ + Global Best $\times 70\%$)
\ENDFOR
\end{algorithmic}
\end{algorithm}


\subsection{CUDA Acceleration}
為了解決 GA-ACO 架構龐大的計算需求，本研究在 GPU 端進行了深度的平行化封裝。
我們設計了 \texttt{ACO\_Environment} 類別，
為 GA 族群中的每一個染色體配置獨立的 GPU 記憶體區塊（包含費洛蒙矩陣、螞蟻狀態與隨機數種子）
以及專屬的並行通道（\texttt{cudaStream\_t}）。
在 CPU 端，利用 C++ 的 \texttt{std::thread} 同時喚醒所有染色體的評估程序；
而在 GPU 端，各染色體的 ACO 模擬被分配至不同的 Stream 異步執行（Asynchronous Execution）。
此設計確保了十組甚至數十組 ACO 演算法能夠在 GPU 內部「真正同時」運行，
且互不干擾記憶體，徹底解決了 Race Condition 的問題，並將整體運算時間縮減至單執行緒的數分之一。

\begin{algorithm}[H]
\caption{Multi-Stream CUDA Execution Architecture}
\begin{algorithmic}[1]
\STATE Array \texttt{envs}[$N$] of \texttt{ACO\_Environment} (each with unique \texttt{cudaStream\_t})
\STATE Array \texttt{threads}[$N$]
\FOR{$i = 1$ to $N$ in CPU}
    \STATE \texttt{threads}[$i$] $\leftarrow$ spawn thread to run \texttt{envs}[$i$].run\_aco($P_i$)
\ENDFOR
\STATE // Inside \texttt{run\_aco} function for environment $i$:
\STATE \texttt{cudaMemcpyAsync}(..., \texttt{envs}[$i$].stream)
\STATE \texttt{ant\_movement\_kernel}$<<<blocks, threads, 0, \text{envs}[i].stream>>>$(...)
\STATE \texttt{cudaMemcpyAsync}(..., \texttt{envs}[$i$].stream)
\STATE \texttt{cudaStreamSynchronize}(\texttt{envs}[$i$].stream)
\STATE Join all \texttt{threads}
\end{algorithmic}
\end{algorithm}


\subsection{Visualizer}
為便於分析與展示演算法成果，本研究開發了整合 OpenCV 與 Gnuplot 的視覺化模組。
在 ACO 的尋優過程中，OpenCV 負責將二維網格地圖渲染為圖形介面，
即時顯示各管線的費洛蒙分佈（以不同顏色的透明度疊加）以及全域最佳路徑的骨架中心線。
當單次 ACO 執行完畢時，系統會自動將最終成果截圖保存。\\
同時，在 GA 演化的巨觀層面，系統利用背景執行續呼叫 Gnuplot，
並將每一代的最佳適應度（Best Score）與中位數適應度（Median Score）即時寫入 Log 檔進行繪圖。
此機制提供了一目了然的演化收斂折線圖，協助開發者即時監控參數優化狀態，
並確保了無頭模式（Headless）下執行與輸出的穩定性。

\begin{algorithm}[H]
\caption{Background Visualization Workflow}
\begin{algorithmic}[1]
\STATE Initialize headless Gnuplot process via \texttt{popen}
\FOR{each GA Generation}
    \STATE Calculate $Best\_Score$ and $Median\_Score$
    \STATE Append scores to \texttt{ga\_log.txt}
    \STATE Send \texttt{plot} commands to Gnuplot pipe
    \STATE Save Gnuplot output as \texttt{ga\_progress.png}
    \STATE OpenCV \texttt{imread} and \texttt{imshow} the chart
    \STATE Render best ACO simulation to OpenCV Mat
    \IF{generation == final}
        \STATE \texttt{imwrite} final ACO path grid to file
    \ENDIF
\ENDFOR
\end{algorithmic}
\end{algorithm}

%%%%%%%%%%%%%%%%%%%%%%%%%%%%%%%%%%%%%%%%%%%%%%%%%%%

\section{Experiment}

\clearpage

\section{Conclusion}

\begin{thebibliography}{00}
\bibitem{b1} W. Jian, Z. Xiaodong, and H. Yisha, ``A novel heuristic approach for 3D pipe routing in complex environments,'' \textit{IEEE Access}, vol. 6, pp. 49764-49775, 2018.
\bibitem{b2} 陳俊仁、翁世光, ``Pipe Routing Method Based on Ant Colony Algorithm'' \textit{國立臺灣海洋大學資訊工程學系碩士學位論文}, 2023.
\bibitem{b3} E. W. Dijkstra, ``A note on two problems in connexion with graphs,'' \textit{Numerische Mathematik}, vol. 1, pp. 269--271, 1959.
\bibitem{b4} P. E. Hart, N. J. Nilsson, and B. Raphael, ``A formal basis for the heuristic determination of minimum cost paths,'' \textit{IEEE Transactions on Systems Science and Cybernetics}, vol. 4, no. 2, pp. 100--107, 1968.
\bibitem{b5} D. E. Goldberg, \textit{Genetic Algorithms in Search, Optimization, and Machine Learning}. Reading, MA: Addison-Wesley, 1989.
\bibitem{b6} M. Dorigo, V. Maniezzo, and A. Colorni, ``Ant system: optimization by a colony of cooperating agents,'' \textit{IEEE Transactions on Systems, Man, and Cybernetics, Part B (Cybernetics)}, vol. 26, no. 1, pp. 29--41, Feb. 1996.
\bibitem{b7} J. M. Cecilia, J. M. García, A. Nisbet, D. Talia, and M. Amos, ``Enhancing data parallelism for Ant Colony Optimization on GPUs,'' \textit{Journal of Parallel and Distributed Computing}, vol. 73, no. 1, pp. 42--51, 2013.
\end{thebibliography}

\end{document}